\section{Couplage réversible et transport non stationnaire}
\label{nuclear:10}
\subsection{Extension réversible du bilan entre deux formes}

Soient \(B,C:V\to V'\) des isomorphismes qui conservent une forme non dégénérée \(b\) entre deux réalisations :



\[
B^*b'B=C^*b'C=b.
\]



La forme peut être une métrique positive ou une forme alternée. Dans les espaces réels, \(*\) désigne la transposition par rapport aux coordonnées choisies ; l’égalité exprime le transport de la forme entre les réalisations spécifiées. Dans les espaces de Hilbert, les opérateurs et leurs inverses sont bornés.


Avec \(s=\sqrt8\) et \(Q_V=\frac13\left(\begin{smallmatrix}sI&-I\\I&sI\end{smallmatrix}\right)\), nous définissons



\begin{equation}
\boxed{
\mathcal U(B,C)=Q_{V'}^*\operatorname{diag}(B,C)Q_V
=\frac19\begin{pmatrix}8B+C&\sqrt8(C-B)\\\sqrt8(C-B)&B+8C\end{pmatrix}.}
\label{nuc10:eq:1}
\end{equation}



Nous écrivons \(T=(8B+C)/9\), \(D=\sqrt8(C-B)/9\) et \(R=(B+8C)/9\). La mise à jour complète est



\begin{equation}
\binom{x'}{m'}=\begin{pmatrix}T&D\\D&R\end{pmatrix}\binom{x}{m}.
\label{nuc10:eq:2}
\end{equation}



\begin{theorem}
L’application \eqref{nuc10:eq:1} est réversible et conserve la forme \(b\oplus b\) :



\begin{equation}
\mathcal U(B,C)^*(b'\oplus b')\mathcal U(B,C)=b\oplus b,
\quad\mathcal U(B,C)^{-1}=\mathcal U(B^{-1},C^{-1}).
\label{nuc10:eq:3}
\end{equation}
\end{theorem}



\begin{proof}
On a \(Q^*Q=I\) par \(8+1=9\). Ses blocs scalaires font que \(Q\) préserve \(b\oplus b\) pour toute forme \(b\). L’application diagonale conserve les formes par hypothèse. La composition prouve \eqref{nuc10:eq:3}, et l’inversion des trois facteurs donne l’inverse déclaré.
\end{proof}

Pour une mémoire entrante \(m=0\), on retrouve le bilan \(b=T^*b'T+D^*b'D\). La seconde colonne détermine la réponse à une mémoire entrante arbitraire. La conservation des normes positives et la conservation de l’aire orientée sont les spécialisations respectives de \eqref{nuc10:eq:3}.

\subsubsection{Réalisation collective dans les neuf copies
}

Dans \(V^9\), soit \(W_E=\operatorname{diag}(I,\ldots,I,E)\), avec \(E\) en neuvième position. Soient \(A(a,b)=(a/\sqrt8,\allowbreak\ldots,\allowbreak a/\sqrt8,b)\) et \(L=AQ\). Alors \(L(x,0)=Jx\) et




\begin{equation}
W_E L=L\mathcal U(I,E).
\label{nuc10:eq:4}
\end{equation}



\begin{proof}
\(W_EA=A\operatorname{diag}(I,E)\) ; multiplier par \(Q\) et utiliser \(AQ=L\) donne \eqref{nuc10:eq:4}. L’orthogonalité de la moyenne des huit premières positions avec leurs sept différences montre en outre que le complément reste fixe sous \(W_E\).
\end{proof}

Le pas cyclique initial est \(S_E=P_{\rm cyc}W_E\). Par conséquent, \(S_EL=(P_{\rm cyc}L)U(I,E)\), avec des systèmes de coordonnées d’entrée et de sortie distincts. La réduction collective conserve le motif correspondant : \(S_E^9=I_9\otimes E\), tandis que \(W_E^9=\operatorname{diag}(I_8,E^9)\). Lorsqu’un lecteur utilise les neuf étiquettes individuelles, celles-ci sont conservées dans leur représentation initiale.

\subsection{Composition exacte et réentrée de mémoire
}

Pour \(B_1,C_1:V_0\to V_1\) et \(B_2,C_2:V_1\to V_2\),



\begin{equation}
\boxed{\mathcal U(B_2,C_2)\mathcal U(B_1,C_1)
=\mathcal U(B_2B_1,C_2C_1).}
\label{nuc10:eq:5}
\end{equation}



\begin{proof}
Dans \eqref{nuc10:eq:1}, \(Q_{V_1}Q_{V_1}^*\) s’annulent, et les deux blocs diagonaux se multiplient dans l’ordre indiqué. La récurrence étend la formule à tout nombre fini de transports, en conservant leur ordre.

\end{proof}

Le bloc de lecture contient l’identité



\begin{equation}
\boxed{T_2T_1+D_2D_1=\frac{8B_2B_1+C_2C_1}{9}.}
\label{nuc10:eq:6}
\end{equation}



À partir de \((x,0)\), la première étape produit \((T_1x,D_1x)\). La seconde lecture est \(T_2T_1x+D_2D_1x\). Le second terme est exactement la contribution de la mémoire réintroduite. En stockant \(D_1x\) hors du couplage suivant et en poursuivant seulement avec \(T_1x\), on obtient le procédé à registres séparés, dont la lecture est \(T_2T_1x\). L’égalité \eqref{nuc10:eq:6} précise la différence entre les deux procédés.

La réalisation \(U\) conserve l’état du couplage et les produits ordonnés \(B_N\cdots B_1\), \(C_N\cdots C_1\). La conservation de tous les préfixes inclut en outre le registre des émissions de HMT, qui maintient chaque facteur avec son ordre.

\subsubsection{Équation covariante d’incrément
}

Pour le registre séparé \(x_{n+1}=T_nx_n\), \(m_n=D_nx_n\), on a




\begin{equation}
x_{n+1}-B_nx_n=m_n/\sqrt8.
\label{nuc10:eq:7}
\end{equation}



Avec \(G_0=I\) et \(G_{n+1}=B_nG_n\), le télescopage donne




\begin{equation}
x_0=G_N^{-1}x_N-\frac1{\sqrt8}\sum_{n<N}G_{n+1}^{-1}m_n.
\label{nuc10:eq:8}
\end{equation}



La mémoire est l’incrément par rapport au transport entre formes \(B_n\), avec la normalisation déclarée. L’égalité \eqref{nuc10:eq:8} est une reconstruction finie algébrique ; son passage à une série infinie requiert la convergence. Le théorème suivant construit l’inverse stable au moyen de la limite d’opérateurs positifs.


\subsection{Mémoire d’une suite non stationnaire}

Dans une réalisation de Hilbert \(H\), soient \(C_n\) des unitaires, avec leur ordre de composition déclaré. Nous écrivons



\[
T_n=(8I+C_n)/9,\quad D_n=\sqrt8(C_n-I)/9,
\quad P_0=I,\quad P_{n+1}=T_nP_n,\quad A_N=P_N^*P_N.
\]



L’indice \(N\) compte les compressions avec stockage séparé, à la différence du transport complet avec réentrée. Pour des métriques variables avec \(B_n\) des isomorphismes isométriques, la conjugaison par \(G_n\) transporte les transports vers un espace de Hilbert fixe et produit ce même cas, avec \(\widetilde C_n=G_{n+1}^{-1}C_nG_n\). L’assertion utilise les métriques des fibres et les isométries déclarées entre elles.

\begin{theorem}[Terminal positif et récupération non stationnaire]
\label{nuclear:terminal-no-estacionario}
Il existe un opérateur positif \(0\leq A_\infty\leq I\), limite forte de \(A_N\), et



\begin{equation}
\boxed{I=A_\infty+\sum_{n\ge0}P_n^*D_n^*D_nP_n.}
\label{nuc10:eq:9}
\end{equation}



L'application



\begin{equation}
\boxed{\mathcal Vv=(A_\infty^{1/2}v,(D_nP_nv)_{n\ge0})}
\label{nuc10:eq:10}
\end{equation}



est isométrique et possède un inverse à gauche stable




\begin{equation}
\mathcal V^*(a,(m_n))=A_\infty^{1/2}a+
\sum_{n\ge0}P_n^*D_n^*m_n,\qquad\|\mathcal V^*\|=1.
\label{nuc10:eq:11}
\end{equation}



\end{theorem}
\begin{proof}
L’identité locale \(T_n^*T_n+D_n^*D_n=I\) donne
\(A_N-A_{N+1}=P_N^*D_N^*D_NP_N\geq0\). Les formes quadratiques de
la suite positive décroissante convergent vers celle d’un opérateur positif
\(A_\infty\). Pour \(B_N=A_N-A_\infty\), l’inégalité
\(0\leq B_N^2\leq B_N\) transforme cette convergence en convergence forte.
Le télescopage fini et la limite prouvent \eqref{nuc10:eq:9}. L’évaluation sur \(v\)
prouve \eqref{nuc10:eq:10}. Son adjoint est \eqref{nuc10:eq:11} ; la série converge en norme parce que
les troncatures d’un registre de carré sommable convergent et que l’adjoint est borné.
\end{proof}

La reconstruction avec le terminal limite et \(N\) registres a pour erreur
exacte \((A_N-A_\infty)v\). Le vecteur \(A_\infty^{1/2}v\) est la
réalisation positive canonique du terminal de norme. Son existence provient
de la limite de Gram, indépendamment de la convergence vectorielle de \(P_Nv\).


\subsection{Trois propriétés différentes : épuisement, distinction et stabilité}

Soit \(Mv=(D_nP_nv)_n\). De \eqref{nuc10:eq:9}, \(M^*M=I-A_\infty\). Par conséquent :

\begin{enumerate}
\item Toute la norme de \(v\) passe au registre exactement lorsque \(A_\infty v=0\).

\item Le registre distingue chaque état exactement lorsque \(\ker(I-A_\infty)=\{0\}\).
\item La récupération à partir de la mémoire seule possède un inverse borné exactement lorsque
\(I-A_\infty\geq\varepsilon I\) pour un certain \(\varepsilon>0\).
Dans ce cas, un inverse à gauche est \((I-A_\infty)^{-1}M^*\).

\end{enumerate}

De plus,



\begin{equation}
\boxed{\ker M=\bigcap_{n\ge0}\operatorname{Fix}(C_n).}
\label{nuc10:eq:12}
\end{equation}



\begin{proof}
Si \(Mv=0\), alors \(D_0v=0\), d’où \(C_0v=v\)
et \(P_1v=v\). Par récurrence, \(D_nP_nv=0\) donne \(C_nv=v\) et
\(P_{n+1}v=v\). La réciproque est immédiate. Les trois équivalences se
déduisent du Gram \(M^*M\) et de la caractérisation des opérateurs
bornés inférieurement.
\end{proof}

\subsubsection{Exemple exact qui sépare les trois propriétés}

Dans l’exemple algébrique \(C_0=-I\) et \(C_n=I\) pour \(n\geq1\),



\begin{equation}
A_\infty=\frac{49}{81}I,\qquad M^*M=\frac{32}{81}I,
\qquad m_0=-\frac{2\sqrt8}{9}v,
\qquad v=-\frac9{2\sqrt8}m_0.
\label{nuc10:eq:13}
\end{equation}



La mémoire contient \(32/81\) du carré de la norme et permet de récupérer le vecteur complet. Le terminal retient \(49/81\) et est un opérateur positif distinct d’un projecteur. La récupérabilité dépend du noyau et du conditionnement du lecteur. Lecture et mémoire constituent des composantes corrélées d’une isométrie.


Cet exemple sépare les trois conséquences logiques du théorème. L’application
à une trajectoire physique conserve en outre la sélection et la
provenance effective de ses \(C_n\) au moyen du TPK.

\subsection{Critère quantitatif pour un transfert complet}

Pour \(P_{j,n}=T_{j-1}\cdots T_n\) et \(P_{n,n}=I\), une fenêtre de longueur \(L\) a pour Gram



\[
G_{n,L}=\sum_{j=n}^{n+L-1}P_{j,n}^*D_j^*D_jP_{j,n}
=I-P_{n+L,n}^*P_{n+L,n}.
\]



Si \(G_{n,L}\geq\kappa I\), avec \(\kappa>0\), dans chaque bloc consécutif de longueur \(L\), alors



\begin{equation}
\|P_N\|^2\le(1-\kappa)^{\lfloor N/L\rfloor},\qquad A_\infty=0.
\label{nuc10:eq:14}
\end{equation}



Chaque bloc contracte la norme au carré par \(1-\kappa\) ; la composition et la contractivité des pas restants prouvent \eqref{nuc10:eq:14}.

Un critère suffisant antérieur aux produits est



\begin{equation}
\sum_{j=n}^{n+L-1}D_j^*D_j\ge\alpha I
\quad\Longrightarrow\quad
G_{n,L}\ge\frac{\alpha}{[1+2(L-1)/9]^2}I.
\label{nuc10:eq:15}
\end{equation}



\begin{proof}
Pour \(y_j=P_{j,n}v\), on a
\(y_{j+1}-y_j=D_jy_j/\sqrt8\) et \(\|D_j\|\leq2\sqrt8/9\). C’est pourquoi

\[
\|D_jv\|\leq\|D_jy_j\|+\frac29\sum_{i=n}^{j-1}\|D_iy_i\|.
\]
La matrice triangulaire des sommes partielles est de norme au plus \(L-1\).
Appliquer la norme \(\ell^2\) de la fenêtre et l’hypothèse de Gram prouve \eqref{nuc10:eq:15}.

\end{proof}

L’absence de vecteur fixe commun donne la distinction, tandis que \eqref{nuc10:eq:15} ajoute une séparation uniforme quantitative. Aucune des deux conditions n’est présumée du seul fait de mentionner des symétries : elle s’applique aux transports effectifs de la séquence.


\subsection{Réalisation d’incidence intégrale et changements de régime}

La section~\ref{nuc05:isometria} construit, pour le registre de douze
coordonnées, \(Fk=(\mu(k),\mathcal IP_0k/6)\), avec inverse et
produit scalaire pondéré dans \(Y=\RR\oplus\mathcal I(\mathbf1^\perp)\).
La transformation \(F\oplus F\) envoie le couplage complet sur



\begin{equation}
\mathcal U_Y=(F'\oplus F')\mathcal U(B,C)(F^{-1}\oplus F^{-1})
=\mathcal U(F'BF^{-1},F'CF^{-1}).
\label{nuc10:eq:16}
\end{equation}



Cette identité conserve la seconde colonne, la réentrée et la composition.
Pour des formes variables, on transporte également la forme au moyen de \(F\).
Dans la réalisation réelle des douze fibres \(A_2\), on utilise
\(F\otimes\operatorname{id}_{V_{A_2}}\), avec
\(V_{A_2}=A_2\otimes_{\ZZ}\RR\). Son domaine est l’espace réel des
fibres ; la préservation du réseau entier correspond au transport
réticulaire déjà démontré. L’application \(F\) retrouve \(K\) et
conserve les autres étiquettes généalogiques dans l’état complet.

La version positive de \eqref{nuc10:eq:9}–\eqref{nuc10:eq:15} utilise des transports unitaires dans les
métriques déclarées. La version à forme alternée de \eqref{nuc10:eq:1}–\eqref{nuc10:eq:8} admet
les trois régimes d’Euler. Avec
\(E=F_{\rm av}^2=\left(
\begin{smallmatrix}1&1\\1&2\end{smallmatrix}\right)\),
on obtient \(\det T=89/81\), \(\det D=-8/81\). Les sommes
orientées finies se télescopent et les terminaux d’aire croissent comme
\((89/81)^N\). La limite positive de Hilbert utilise les hypothèses
unitaires spécifiques du théorème~\ref{nuclear:terminal-no-estacionario}.


Le bilan réunit les formes et les domaines, l’évolution réversible avec réentrée
et la récupération illimitée dans les réalisations positives. Les
applications physiques suivantes conservent les lecteurs, les connexions et les
échelles de cette composition opératorielle.
